Una cooperativa agrícola dispone de tres tipos de cultivos: \textbf{trigo}, \textbf{maíz}  y \textbf{girasol}  
Cada uno requiere una cierta cantidad de agua y superficie, y aporta una \textbf{producción total de biomasa (en toneladas)} que la cooperativa quiere \textbf{maximizar}.

Los datos son los siguientes:

\begin{center}
\begin{tabular}{lccc}
\toprule
\textbf{Cultivo} & \textbf{Agua necesaria (m$^3$/ha)} & \textbf{Biomasa (ton/ha)} \\
\midrule
\textbf{Trigo} ($x_1$) & 300 & 4 \\
\textbf{Maíz} ($x_2$) & 500 & 6 \\
\textbf{Girasol} ($x_3$) & 400 & 5 \\
\bottomrule
\end{tabular}
\end{center}

Cada semana, la cooperativa dispone de \textbf{30\,000 m$^3$ de agua} y \textbf{80 hectáreas de terreno total}.  
Además, por exigencias ecológicas, el cultivo de maíz y girasol juntos debe ocupar \textbf{al menos 40 hectáreas}.

El modelo de Programación Lineal que permite determinar el plan de cultivo que maximiza la producción de biomasa es:

\[
\begin{aligned}
\max \quad & z = 4x_1 + 6x_2 + 5x_3 \\
\text{r. t.} \quad 
& 3x_1 + 5x_2 + 4x_3 \le 300 \quad && \text{(disponibilidad de agua)} \\
& x_1 + x_2 + x_3 \le 80 \quad && \text{(superficie total)} \\
& x_2 + x_3 \ge 40 \quad && \text{(requisito ecológico)} \\
& x_1, x_2, x_3 \ge 0
\end{aligned}
\]

donde $x_1$, $x_2$ y $x_3$ representan, respectivamente, las hectáreas dedicadas a trigo, maíz y girasol,

Se conoce la tabla correspondiente a la solución óptima de dicho problema:

\begin{center}
\begin{tabular}{c|c|ccccccc|}
 & $z$ & $x_1$ & $x_2$ & $x_3$ & $h_1$ & $h_2$ & $h_3$ & $a_3$\\ 
      & $-380$ & $0$ & $0$ & $0$ & $-1$ & $-1$ & $0$ & $0$ \\ 
\hline
$x_1$ & 40  & $1$ & $0$ & $0$ & $0$ & $1$ & $1$ & $-1$\\ 
$x_3$ & 20  & $0$ & $0$ & $1$ & $-1$ & $3$ & $-2$ & $2$\\ 
$x_2$ & 20  & $0$ & $1$ & $0$ & $1$ & $-3$ & $1$ & $-1$\\ 
\hline
\end{tabular}
\end{center}
 
\bigskip
Se pide lo siguiente:
\begin{enumerate}[label=\alph*)]
    \item Interpretar la solución óptima e indicar cuál es el plan de cultivo óptimo de la cooperativa y del cumplimiento de los requisitos descritos.

    % Impacto sobre el beneficio y sobre el plan de producción del cambio de B3 = 1
    \item Cuál sería el impacto sobre el plan de producción y sobre la producción total de biomasa si se incrementase en una unidad el requisito ecológico relativo al mínimo del total de producción de maíz y girasol.

    % Sensibilidad de c_1
    \item Cuál es la producción de masa por hectárea de maíz, por debajo de la cual no sería interesante cultivar maíz

    % Nuevo cultivo
    \item Existe la posibilidad de introducir un cuarto cultivo, \textbf{soja}, que genera biomasa.
    Se sabe que requiere \(\,350\,\text{m}^3/\text{ha}\) de agua y produce 4.5 ton/ha).
    La soja no forma parte del requisito ecológico (esto es, no contribuye al mínimo conjunto de maíz y girasol).
    
    % Restricción de mercado de un producto (Lemke)
    \item Debido a las previsiones meteorológicas para la temporada, la cooperativa ha decidido producir un máximo de 10 hectáreas de maíz, ¿cuál sería el nuevo plan de cultivo óptimo en este caso? 
    
\end{enumerate}
